% ##############################################################################
\section{Conclusion}
\label{sec:conclusion}

This paper described the canary subsystem of the \Noesis{} prototype, which
instantiates the white paper's pluggable capability measurement with
known-answer probes. The design combines a deterministic bank and grader, a
calibration step that measures each candidate against healthy and weak targets,
keyed selection and scheduling with provable balance and spacing, a prober that
sends probes as ordinary user messages through the gateway, and a verdict rule
that flags a seller only when an upper confidence bound falls below the
promise. Exact binomial analysis indicates that, with the prototype defaults of
ten canaries and a promised accuracy of 0.85, the rule rarely flags a compliant
seller and flags a seller with canary accuracy 0.5 in about 83\% of contracts,
while shortfalls near the promise need hundreds of probes.

The evidence is analytical and unit-level. The live calibration is pending, the
database-backed tests were not run in the reviewed environment, and the main
threats to validity, namely recognizable prompts, a single weak reference
model, and strict exact-match grading, have not been measured. The next step is
the reviewed calibration run and an end-to-end fault-injection experiment, which
would replace the assumptions of this paper with measured accuracies and
detection rates.
